% =====================================================================
%  BAB 3 — METODOLOGI PENELITIAN
% =====================================================================

\chapter{Metodologi Penelitian}

\section{Desain Penelitian}
\label{sec:3.1}

Penelitian ini menggunakan pendekatan [eksperimental / deskriptif / kuantitatif / kualitatif / campuran].
[Jelaskan alasan pemilihan desain penelitian ini dan bagaimana desain tersebut menjawab
rumusan masalah yang telah ditetapkan.]

% Contoh gambar alur penelitian
\begin{figure}[H]
  \centering
  \begin{tikzpicture}[node distance=1cm]
    \node[term]    (mulai)   {Mulai};
    \node[proc]    (studi)   [below=of mulai]   {Studi Literatur};
    \node[proc]    (rumus)   [below=of studi]   {Perumusan Masalah};
    \node[proc]    (data)    [below=of rumus]   {Pengumpulan Data};
    \node[proc]    (analisis)[below=of data]    {Analisis \& Perancangan};
    \node[proc]    (impl)    [below=of analisis]{Implementasi};
    \node[proc]    (uji)     [below=of impl]    {Pengujian \& Evaluasi};
    \node[dec]     (valid)   [below=of uji]     {Tujuan \\ tercapai?};
    \node[proc]    (perbaiki)[right=2cm of valid]{Perbaikan};
    \node[proc]    (laporan) [below=of valid]   {Penulisan Laporan};
    \node[term]    (selesai) [below=of laporan] {Selesai};

    \draw[arr] (mulai)    -- (studi);
    \draw[arr] (studi)    -- (rumus);
    \draw[arr] (rumus)    -- (data);
    \draw[arr] (data)     -- (analisis);
    \draw[arr] (analisis) -- (impl);
    \draw[arr] (impl)     -- (uji);
    \draw[arr] (uji)      -- (valid);
    \draw[arr] (valid)    -- node[flbl,right]{Ya}  (laporan);
    \draw[arr] (valid)    -- node[flbl,above]{Tidak}(perbaiki);
    \draw[arr] (perbaiki) |- (impl);
    \draw[arr] (laporan)  -- (selesai);
  \end{tikzpicture}
  \caption{Alur penelitian}
  \label{fig:3.alur}
\end{figure}

\section{Dataset dan Lingkungan Eksperimen}
\label{sec:3.2}

\subsection{Dataset}
\label{sec:3.2.1}

[Deskripsikan dataset yang digunakan: sumber, ukuran, karakteristik distribusi kelas,
format, dan cara mendapatkannya. Jelaskan apakah dataset publik atau dikumpulkan sendiri.]

\begin{table}[H]
  \centering
  \caption{Karakteristik dataset yang digunakan}
  \label{tab:3.dataset}
  \begin{tabular}{lrrl}
    \toprule
    \textbf{Subset} & \textbf{Jumlah sampel} & \textbf{Proporsi (\%)} & \textbf{Keterangan} \\
    \midrule
    {[Kelas A / Benign]} & [X.XXX] & [XX,X] & [Deskripsi singkat] \\
    {[Kelas B / Attack]} & [X.XXX] & [XX,X] & [Deskripsi singkat] \\
    \midrule
    \textbf{Total}     & [X.XXX] & 100,0  & \\
    \bottomrule
  \end{tabular}
\end{table}

\subsection{Lingkungan Eksperimen}
\label{sec:3.2.2}

Eksperimen dilakukan pada lingkungan sebagai berikut:

\begin{table}[H]
  \centering
  \caption{Spesifikasi lingkungan eksperimen}
  \label{tab:3.lingkungan}
  \begin{tabular}{ll}
    \toprule
    \textbf{Komponen} & \textbf{Spesifikasi} \\
    \midrule
    Sistem Operasi   & [Nama OS, Versi] \\
    Prosesor         & [Nama CPU, Jumlah Core, Frekuensi] \\
    Memori (RAM)     & [Kapasitas GB] \\
    Penyimpanan      & [Tipe, Kapasitas] \\
    Bahasa Pemrograman & [Bahasa, Versi] \\
    Framework/Library& [Nama, Versi] \\
    \bottomrule
  \end{tabular}
\end{table}

\section{Metode yang Digunakan}
\label{sec:3.3}

\subsection{[Metode 1]}
\label{sec:3.3.1}

[Uraikan metode/algoritma pertama yang digunakan. Jelaskan cara kerja, parameter kunci,
dan alasan pemilihan metode ini. Sertakan pseudocode jika perlu.]

% Contoh pseudocode
% \begin{lstlisting}[language=Python, caption={Pseudocode [Metode 1]}, label={lst:3.pseudo}]
% def metode(input):
%     # Langkah 1: inisialisasi
%     hasil = []
%     # Langkah 2: proses
%     for item in input:
%         hasil.append(proses(item))
%     return hasil
% \end{lstlisting}

\subsection{[Metode 2]}
\label{sec:3.3.2}

[Uraikan metode kedua.]

\section{Metrik Evaluasi}
\label{sec:3.4}

Kinerja sistem dievaluasi menggunakan metrik-metrik berikut.

\textbf{Presisi} (\emph{Precision}) mengukur proporsi prediksi positif yang benar:
\begin{equation}
  \text{Presisi} = \frac{TP}{TP + FP}
  \label{eq:presisi}
\end{equation}

\textbf{Recall} (\emph{Sensitivitas}) mengukur kemampuan model mendeteksi kasus positif:
\begin{equation}
  \text{Recall} = \frac{TP}{TP + FN}
  \label{eq:recall}
\end{equation}

\textbf{F1-score} merupakan rata-rata harmonis antara presisi dan recall:
\begin{equation}
  F_1 = 2 \cdot \frac{\text{Presisi} \times \text{Recall}}{\text{Presisi} + \text{Recall}}
  \label{eq:f1}
\end{equation}

\noindent
di mana $TP$ = \emph{True Positive}, $FP$ = \emph{False Positive}, $FN$ = \emph{False Negative}.

[Tambahkan metrik lain yang relevan, misalnya latensi, throughput, AUC-ROC, dsb.]

\section{Prosedur Eksperimen}
\label{sec:3.5}

[Jelaskan langkah-langkah eksperimen secara runtut: pembagian data (train/val/test),
metode validasi (k-fold cross-validation, hold-out), skenario ablasi, dsb. Pastikan
eksperimen dapat direproduksi oleh peneliti lain.]

Eksperimen dilaksanakan dalam tahapan berikut:

\begin{enumerate}
  \item \textbf{Praproses data}: [normalisasi, penanganan nilai hilang, encoding, dsb.]
  \item \textbf{Pelatihan model}: [pembagian data, hiperparameter, strategi tuning.]
  \item \textbf{Pengujian}: [skenario pengujian, baseline yang dibandingkan.]
  \item \textbf{Analisis}: [metrik yang dihitung, visualisasi yang dibuat.]
\end{enumerate}
